\section{The dense supporting families}
\label{sec:families}

GEMV, TRSM and the cuSOLVER wrapper round out the library. Each is correct against
the vendor library and benchmarked on the same device, and Table \ref{tab:families}
is generated from whichever of them the committed summary carries rows for.

GEMV is a bandwidth benchmark in disguise. Its arithmetic intensity is fixed near
0.5 FLOP per byte whatever the size, so the warp and vectorized kernels reach both
cuBLAS and the measured streaming bandwidth, while the naive kernel is coalescing
limited and never gets near either. On the roofline it sits on the bandwidth
diagonal, which is the point of including it: the same chart that shows GEMM far
into the compute bound region shows GEMV pinned to the memory bound one.

The v1 table had the warp kernel at 2048 reading 208.6 percent of cuBLAS, which was
not a result but the defect correction A2 names: setup work inside a vendor timed
region. With the wrapper fixed and the sweep re-run at locked clocks, that row
reads 111.2 percent. Four GEMV rows still sit above 100, and at this point they
are believable rather than suspicious: the vectorized and warp kernels beat cuBLAS
SGEMV at 1024 and 2048 square, by 17.7 percent at most, and both fall back under it
at 4096 and 8192, where the problem is large enough that the vendor's own
decomposition wins. A memory bound kernel beating a vendor library on a small
problem is an ordinary outcome, and the way to tell it from a measurement defect
was to fix the defect and look again.

TRSM reduces most of its work to a GEMM style trailing update. The blocked kernel
agrees with cuBLAS to about $10^{-7}$, which is a correctness result and not a
performance one: on speed it reaches 322.8 GFLOP/s at 2048 by 256 against cuBLAS
STRSM's 1308.9, so 24.7 percent, and closing that gap is open work. That percentage
is \emph{lower} than the 31.8 the v1 table reported, and the reason is worth
stating: correction A2 removed a host to device restore copy from inside the timed
region on both sides, and taking setup out of the window made the vendor baseline
faster by more than it made my kernel faster. Fixing a measurement defect is
supposed to be able to move a number in the unflattering direction, and this one
did. The naive column at a time kernel is kept in the table as the honest baseline
rather than deleted for being slow.

The cuSOLVER LU and Cholesky paths are checked by residual, about $5.6 \times
10^{-7}$ and $4.3 \times 10^{-7}$, rather than merely by the absence of an error
code. They are a correctness deliverable and a RAII wrapper, not a performance
claim, and they carry no timed row.

\begin{table}[htbp]
\centering
\begin{tabular}{llrrl}
\toprule
kernel & shape & GFLOP/s & percent of vendor & baseline \\
\midrule
GEMV naive & 1024 sq & 37.9 & 15.8 & cuBLAS SGEMV \\
GEMV vectorized & 1024 sq & 227.2 & 134.7 & cuBLAS SGEMV \\
GEMV warp & 1024 sq & 225.6 & 93.8 & cuBLAS SGEMV \\
GEMV naive & 2048 sq & 82.7 & 15.0 & cuBLAS SGEMV \\
GEMV vectorized & 2048 sq & 586.5 & 101.7 & cuBLAS SGEMV \\
GEMV warp & 2048 sq & 635.5 & 208.6 & cuBLAS SGEMV \\
GEMV naive & 4096 sq & 134.0 & 46.0 & cuBLAS SGEMV \\
GEMV vectorized & 4096 sq & 291.4 & 100.9 & cuBLAS SGEMV \\
GEMV warp & 4096 sq & 291.7 & 101.4 & cuBLAS SGEMV \\
GEMV naive & 8192 sq & 159.5 & 52.2 & cuBLAS SGEMV \\
GEMV vectorized & 8192 sq & 308.1 & 101.6 & cuBLAS SGEMV \\
GEMV warp & 8192 sq & 304.0 & 99.9 & cuBLAS SGEMV \\
TRSM blocked & 512 by 256 & 109.9 & 47.4 & cuBLAS STRSM \\
TRSM naive & 512 by 256 & 8.4 & 3.2 & cuBLAS STRSM \\
TRSM blocked & 1024 by 256 & 215.5 & 37.7 & cuBLAS STRSM \\
TRSM naive & 1024 by 256 & 7.9 & 1.4 & cuBLAS STRSM \\
TRSM blocked & 2048 by 256 & 381.3 & 31.8 & cuBLAS STRSM \\
TRSM naive & 2048 by 256 & 7.8 & 0.7 & cuBLAS STRSM \\
\bottomrule
\end{tabular}

\caption{The supporting families on the RTX 5070, each versus its vendor baseline,
generated from the families the committed summary actually carries.}
\label{tab:families}
\end{table}
